% appendices/appendix-c-glosarium.tex
\chapter{Glosarium}
\label{app:glosarium}

Glosarium ini memuat istilah yang dipakai pada seluruh buku, disusun menurut
abjad istilah Indonesianya. Kolom kedua memuat padanan Inggrisnya, sebab
sebagian besar dokumentasi dan pesan kesalahan yang akan Anda temui memakai
padanan itu.

\begin{longtable}{@{}L{3.9cm} L{3.6cm} L{7.2cm}@{}}
\caption{Glosarium istilah Indonesia--Inggris}
\label{tab:glosarium} \\
\toprule
\headrow \textbf{Istilah} & \textbf{Padanan Inggris} & \textbf{Keterangan singkat} \\
\midrule
\endfirsthead
\toprule
\headrow \textbf{Istilah} & \textbf{Padanan Inggris} & \textbf{Keterangan singkat} \\
\midrule
\endhead
\bottomrule
\endlastfoot

Algoritma pencarian jalur & \istilah{Pathfinding algorithm} &
Algoritma yang mencari rangkaian petak dari satu titik ke titik lain, misalnya
BFS, Dijkstra, dan A*. \\ \addlinespace

Animasi & \istilah{Animation} &
Rangkaian gambar yang diputar berurutan sehingga tampak bergerak. \\
\addlinespace

Antarmuka pengguna & \istilah{User interface} (UI) &
Bagian layar yang bukan bagian dari dunia permainan, misalnya skor, nyawa, dan
tombol. \\ \addlinespace

Antrean prioritas & \istilah{Priority queue} &
Struktur data yang mengeluarkan isinya menurut prioritas, bukan menurut urutan
masuk. A* memakainya untuk memilih petak berikutnya. \\ \addlinespace

Aset & \istilah{Asset} &
Berkas yang dipakai permainan: gambar, bunyi, huruf, dan animasi. \\
\addlinespace

Batang (badan) & \istilah{Body} &
Benda yang diperlakukan sebagai satu kesatuan gerak, dengan posisi, kecepatan,
dan bentuk. \\ \addlinespace

Batas benda & \istilah{Collider} &
Bentuk yang dipakai mesin untuk memeriksa tabrakan, tidak selalu sama dengan
bentuk yang terlihat pemain. \\ \addlinespace

Bingkai & \istilah{Frame} &
Satu kali gambar lengkap ditampilkan. Pada 60 bingkai per detik, satu bingkai
berlangsung sekitar 16,7 milidetik. \\ \addlinespace

Bunyi latar & \istilah{Background music} (BGM) &
Bunyi panjang yang berjalan terus selama permainan, dan tidak dipanggil
menurut kejadian. \\ \addlinespace

Bug & \istilah{Bug} &
Keadaan yang tidak seharusnya terjadi, dan dapat diulang dengan urutan
tindakan tertentu. \\ \addlinespace

Build & \istilah{Build} &
Berkas permainan yang sudah dibungkus sehingga dapat dijalankan pada komputer
lain. \\ \addlinespace

Celah & \istilah{Gap} &
Jarak kosong yang harus dilompati pemain di antara dua pijakan. \\
\addlinespace

Delta time & \istilah{Delta time} &
Selisih waktu antara dua bingkai, dipakai agar gerak tidak bergantung pada
kecepatan komputer. \\ \addlinespace

Desibel & \istilah{Decibel} (dB) &
Satuan kekuatan bunyi yang naik secara berlipat. Turun 6 dB berarti kekuatannya
menjadi kira-kira separuh. \\ \addlinespace

Deteksi tabrakan & \istilah{Collision detection} &
Pemeriksaan apakah dua benda saling bersentuhan pada satu bingkai. \\
\addlinespace

Efek suara & \istilah{Sound effect} (SFX) &
Bunyi pendek yang dipanggil menurut peristiwa, dan dapat berbunyi bersamaan
dengan efek suara lain. \\ \addlinespace

Gaya berat & \istilah{Gravity} &
Percepatan tetap yang menarik benda ke bawah. Pada Bab~VI dipakai 1800 piksel
per detik kuadrat. \\ \addlinespace

Graf & \istilah{Graph} &
Kumpulan simpul dan hubungan di antaranya. Peta yang diubah menjadi pijakan
dan lompatan adalah sebuah graf. \\ \addlinespace

Heuristik & \istilah{Heuristic} &
Perkiraan jarak ke tujuan yang dipakai A* untuk mengarahkan pencarian. \\
\addlinespace

Hierarki & \istilah{Hierarchy} &
Susunan induk dan anak pada \istilah{scene}, misalnya \istilah{Canvas} yang
membawahi teks dan ikon. \\ \addlinespace

Impor & \istilah{Import} &
Memasukkan berkas aset ke dalam \istilah{project}, sehingga berkas itu diolah
dan disimpan dalam bentuk yang siap dipakai. \\ \addlinespace

Indikator & \istilah{Indicator} &
Perilaku yang dapat diamati, yang dipakai untuk menyatakan bahwa sebuah
Sub-CPMK telah dikuasai. \\ \addlinespace

Jangkauan & \istilah{Reach} &
Jarak terjauh yang dapat ditempuh pemain dalam satu lompatan, dihitung dari
kecepatan lompatan dan gravitasi. \\ \addlinespace

Kecerdasan buatan & \istilah{Artificial intelligence} (AI) &
Pada buku ini, aturan yang menentukan perilaku NPC, bukan pembelajaran mesin. \\
\addlinespace

Kecepatan bingkai & \istilah{Frame rate} &
Banyaknya bingkai yang ditampilkan dalam satu detik, dinyatakan dalam
\istilah{fps}. \\ \addlinespace

Keadaan & \istilah{State} &
Keadaan yang ditempati sebuah benda atau sistem pada satu saat, misalnya diam,
berlari, atau melompat. \\ \addlinespace

Kendali pemain & \istilah{Player controller} &
Skrip yang membaca masukan pemain dan mengubahnya menjadi gerak benda. \\
\addlinespace

Kode sumber & \istilah{Source code} &
Berkas teks berisi program sebelum dijalankan atau dibangun. \\
\addlinespace

Komponen & \istilah{Component} &
Bagian yang ditempelkan pada \istilah{GameObject} untuk memberi kemampuan,
misalnya \istilah{Rigidbody} atau \istilah{Animator}. \\ \addlinespace

Lembar \istilah{playtest} & \istilah{Playtest sheet} &
Lembar yang dipakai penguji untuk mencatat waktu, tempat, dan keadaan yang
tidak seharusnya terjadi. \\ \addlinespace

Level & \istilah{Level} &
Satu ruang permainan beserta seluruh isinya, dari titik awal sampai tujuan. \\
\addlinespace

Log \istilah{bug} & \istilah{Bug log} &
Daftar temuan beserta urutan yang menyebabkannya, beratnya, dan status
penanganannya. \\ \addlinespace

Masukan & \istilah{Input} &
Apa pun yang dikirim pemain ke program: tombol, tetikus, atau sentuhan. \\
\addlinespace

Mekanik & \istilah{Mechanic} &
Aturan yang menentukan cara pemain memengaruhi permainan, misalnya melompat
dan mengumpulkan benda. \\ \addlinespace

Mesin keadaan & \istilah{State machine} (FSM) &
Susunan keadaan beserta aturan perpindahan di antaranya. \\
\addlinespace

Mesin permainan & \istilah{Game engine} &
Perangkat lunak yang mengurus tampilan, fisika, bunyi, dan aset, misalnya
Unity. \\ \addlinespace

Nada & \istilah{Pitch} &
Tinggi rendahnya bunyi. Menggeser nada sedikit membuat bunyi yang berulang
tidak terdengar seperti mesin. \\ \addlinespace

NPC & \istilah{Non-player character} &
Benda yang bertindak menurut aturan program, bukan menurut masukan pemain. \\
\addlinespace

Objek permainan & \istilah{GameObject} &
Satuan dasar pada Unity: benda yang memiliki letak, dan memperoleh kemampuan
melalui komponen. \\ \addlinespace

Partikel & \istilah{Particle} &
Gambar kecil yang muncul, bergerak, lalu menghilang, dipakai sebagai pasangan
yang terlihat bagi bunyi. \\ \addlinespace

Peristiwa & \istilah{Event} &
Kejadian yang dibaca satu kali, bukan keadaan yang berlangsung, misalnya
tombol ditekan atau benda disentuh. \\ \addlinespace

Petak & \istilah{Tile} &
Satuan kotak pada peta dua dimensi. Jangkauan lompatan pada buku ini dinyatakan
dalam petak. \\ \addlinespace

Pijakan & \istilah{Platform} &
Petak kosong yang di bawahnya berupa dinding, sehingga pemain dapat berdiri di
sana. \\ \addlinespace

\istilah{Pixels Per Unit} & \istilah{Pixels per unit} &
Banyaknya piksel gambar yang mewakili satu satuan pada \istilah{scene}.
Menentukan ukuran gambar di dalam level. \\ \addlinespace

\istilah{Playtest} & \istilah{Playtest} &
Sesi ketika orang lain memainkan prototipe Anda, sementara Anda mengamati dan
mencatat. \\ \addlinespace

Portofolio & \istilah{Portfolio} &
Kumpulan artefak pekerjaan beserta keterangannya, disusun sebagai bukti
kemampuan. \\ \addlinespace

Prototipe & \istilah{Prototype} &
Versi awal yang sudah dapat dimainkan, dipakai untuk menguji rancangan sebelum
dikerjakan sepenuhnya. \\ \addlinespace

\istilah{Rigidbody} & \istilah{Rigid body} &
Komponen Unity yang menyerahkan gerak benda kepada mesin fisika, sehingga
gravitasi dan tabrakan dihitung oleh mesin. \\ \addlinespace

Simulasi & \istilah{Simulation} &
Menjalankan permainan tanpa pemain, berulang kali, untuk memperoleh gambaran
tentang sebagian keadaannya. \\ \addlinespace

Skrip & \istilah{Script} &
Berkas berisi kode yang ditempelkan pada \istilah{GameObject} sebagai komponen
pada Unity. \\ \addlinespace

\istilah{Soft-lock} & \istilah{Soft-lock} &
Keadaan ketika pemain tidak dapat melanjutkan dan tidak dapat kembali,
walaupun programnya masih berjalan dan tidak mencetak kesalahan. \\
\addlinespace

\istilah{Spatial Blend} & \istilah{Spatial blend} &
Pengaturan pada \code{AudioSource} yang menentukan seberapa besar letak bunyi
berpengaruh: nol untuk musik latar, satu untuk bunyi di dalam level. \\
\addlinespace

\istilah{Sprite} & \istilah{Sprite} &
Gambar dua dimensi yang ditempatkan di dalam \istilah{scene}. \\
\addlinespace

\istilah{Sprite sheet} & \istilah{Sprite sheet} &
Satu berkas gambar yang memuat beberapa bingkai, dipotong menjadi bingkai
terpisah. \\ \addlinespace

Sub-CPMK & \istilah{Sub-learning outcome} &
Turunan CPMK pada satu pokok bahasan, beserta indikator ketercapaiannya. \\
\addlinespace

Tabrakan & \istilah{Collision} &
Persentuhan dua benda yang mengubah geraknya, berbeda dari
\istilah{trigger} yang hanya mendeteksi sentuhan. \\ \addlinespace

Tembus & \istilah{Tunneling} &
Keadaan ketika benda berpindah melewati benda lain dalam satu langkah, karena
langkahnya terlalu besar. \\ \addlinespace

Tingkat (audio) & \istilah{Mixer group} &
Kelompok pada \istilah{AudioMixer} yang mengatur kekuatan sekumpulan bunyi
sekaligus. \\ \addlinespace

Transform & \istilah{Transform} &
Komponen yang menyimpan letak, rotasi, dan ukuran sebuah
\istilah{GameObject}. \\ \addlinespace

\istilah{Trigger} & \istilah{Trigger} &
Bentuk batas benda yang mendeteksi sentuhan tanpa menghentikan gerak, dipakai
untuk pintu, benda yang dapat diambil, dan daerah tujuan. \\
\addlinespace

Umpan balik & \istilah{Feedback} &
Tanggapan yang diterima pemain atas tindakannya, melalui tampilan, bunyi, atau
gerak. \\ \addlinespace

Vektor & \istilah{Vector} &
Besaran yang memiliki besar dan arah, dipakai untuk menyatakan posisi dan
kecepatan. \\ \addlinespace

Waktu bingkai & \istilah{Frame time} &
Lama satu bingkai berlangsung. Kebalikan dari kecepatan bingkai. \\
\addlinespace

\end{longtable}
